% !TeX root = ../main.tex
% =========================================================================
% 3. JUSTIFICACIÓN (técnica, social, aporte, factibilidad, restricciones) + alcance
% =========================================================================
\section{Justificación}

La presente investigación se justifica por la imperiosa necesidad de transformar el proceso operativo de cotización y asesoramiento técnico de materiales eléctricos en las distribuidoras comerciales que gestionan sus ventas a través de mensajería instantánea. En el mercado ecuatoriano de suministros eléctricos, la inmediatez y la certeza técnica determinan el cierre comercial; no obstante, el personal de ventas invierte una fracción considerable de su jornada laboral consultando hojas de cálculo, archivos PDF desactualizados y listas impresas para responder preguntas repetitivas sobre marcas, calibres, amperajes y precios. La incorporación de un agente inteligente de ventas comerciales capaz de guiar al cliente en el embudo técnico, interpretar requerimientos conversacionales y emitir cotizaciones fácticas en tiempo real con datos verificables confiere una ventaja competitiva directa y mitiga el riesgo de pérdida de clientes por demoras en la atención.

Desde una perspectiva técnica y científica, la integración de un modelo de lenguaje generativo con una capa de recuperación local estructurada resuelve de manera determinante el fenómeno de las alucinaciones probabilísticas y la inconsistencia fáctica en modelos fundacionales \parencite{ji2023hallucination,huang2023survey,tonmoy2024comprehensive,lewis2020retrieval}. En el campo de los suministros eléctricos, las alucinaciones no constituyen meros descuidos estilísticos; la invención de un calibre de conductor indebido, la alteración de la curva de disparo de un disyuntor termomagnético o la cita de un precio desactualizado acarrean riesgos severos de sobrecarga eléctrica, cortocircuito, siniestros materiales o pérdidas financieras. Si bien los modelos fundacionales como Gemini poseen una alta capacidad de comprensión sintáctica y semántica en lenguaje natural, carecen de acceso autónomo a inventarios específicos y dinámicos. El empleo del motor relacional embebido SQLite con registro de escritura anticipada (WAL) permite realizar consultas deterministas de alta velocidad directamente en el espacio de memoria del proceso, inyectando los datos fácticos verificados (\textit{ground truth}) dentro del contexto de inferencia. Esto erradica las discrepancias de catálogo y evita el despliegue de bases de datos vectoriales pesadas que demandan hardware costoso \parencite{gaffney2022sqlite}. Asimismo, la implementación del servicio en .NET 8 con patrones asíncronos bajo el paradigma de microservicios ligeros orientados a la nube garantiza la capacidad de atender múltiples sesiones simultáneas bajo consumo mínimo de memoria, en concordancia con los principios de desarrollo cloud-native y software asistido por inteligencia artificial promovidos en la disciplina \parencite{kratzke2017understanding,newman2021building,taibi2020patterns,Llerena-Izquierdo2024}.

En el ámbito social, productivo y de política pública, el proyecto se alinea de forma explícita con los objetivos del Plan Nacional de Desarrollo de la República del Ecuador, en sus ejes orientados a la transformación digital, la modernización de los encadenamientos productivos y el fortalecimiento de la competitividad de las micro, pequeñas y medianas empresas (MIPYMES). La distribuidora Power Electrimk, ubicada en la ciudad de Guayaquil, se beneficiará de una operatividad continua 24/7, canalizando prospectos calificados hacia sus ejecutivos de venta. Al mismo tiempo, los maestros eléctricos, contratistas y usuarios finales obtendrán especificaciones técnicas normalizadas y precios referenciales transparentes sin barreras geográficas ni demoras de atención.

En conformidad con los criterios de diseño en ingeniería (Criterio ABET 1), la solución se encuentra sujeta a cuatro restricciones formales de diseño:
\begin{enumerate}
    \item \textbf{Restricción de latencia operativa:} El tiempo total de procesamiento y respuesta de extremo a extremo no deberá superar los 4.0 segundos por interacción, garantizando una experiencia fluida y evitando cancelaciones o reintentos en los webhooks de WhatsApp Cloud API.
    \item \textbf{Restricción de hardware y costo de infraestructura:} El sistema deberá operar de manera estable en una instancia de entrada en la nube de bajos recursos (computación contenida a 1 vCPU y 1 GB de memoria RAM, equivalente a AWS EC2 t2.micro), manteniendo los costos operativos de infraestructura por debajo de 15.00 USD mensuales.
    \item \textbf{Restricción de cumplimiento normativo y privacidad:} El manejo de identificadores de usuario, números telefónicos y registros de conversaciones deberá cumplir estrictamente con los mandatos de la Ley Orgánica de Protección de Datos Personales (LOPDP) de la República del Ecuador, empleando transporte seguro bajo TLS 1.3 y disociación de datos de contacto.
    \item \textbf{Restricción de exactitud e integridad de catálogo:} Se impone una tolerancia de error nula (0\%), supeditando cualquier emisión de cotización a la coincidencia determinista con los registros oficiales de la base de datos local.
\end{enumerate}

El alcance del proyecto de titulación comprende la recopilación y estructuración de los catálogos multimarca de la distribuidora, el diseño del esquema relacional local, el desarrollo del servicio backend en .NET 8 con integración a WhatsApp Cloud API y al modelo generativo Gemini, el despliegue productivo en una instancia Linux de AWS con proxy inverso Nginx y certificados TLS, y la evaluación experimental cuantitativa de latencia, exactitud y consumo de recursos. Se delimita la solución a la asesoría técnica de compatibilidad, emisión de cotizaciones referenciales formales y captura estructurada de prospectos. Asimismo, la arquitectura de software se diseña de manera modular y extensible bajo principios SOLID, sentando las bases contractuales para una segunda fase tecnológica post-titulación que incorporará reconocimiento multimodal de piezas por imágenes fotográficas con Gemini Flash Vision, seguimiento comercial automatizado (\textit{follow-up}), derivación asistida a asesores humanos mediante el patrón \textit{State} y sincronización bidireccional con sistemas CRM. Quedan explícitamente excluidos del alcance del anteproyecto la recaudación de pagos directos mediante pasarelas financieras, la emisión de facturación electrónica reglamentada por el SRI y la conexión nativa con sistemas ERP legados, salvaguardando la viabilidad de ejecución dentro del límite de 400 horas académicas.
